\documentclass[10pt,a4paper]{amsart}
\usepackage[margin=19mm]{geometry}
\usepackage{amsmath,amssymb,mathtools}
\usepackage[scr=rsfs,cal=euler]{mathalfa}
\usepackage{xfrac}
\usepackage[unicode,hidelinks,bookmarks=false]{hyperref}
\newcommand{\ie}{i.e.\ }
\newcommand{\cf}{cf.\ }
\newcommand{\resp}{respectively\ }
\newcommand{\Subsection}{\subsection}
\providecommand{\nobreakdash}{\nobreak\hbox{-}}
\setlength{\emergencystretch}{2em}
\multlinegap0pt
\usepackage{bm}
\usepackage{bbm}
\usepackage{dsfont}
\usepackage{xspace}
\usepackage{cancel}
\usepackage{mathtools}
\usepackage{amsthm}
\makeatletter
\@ifundefined{theorem}{\newtheorem{theorem}{Theorem}}{}
\@ifundefined{lemma}{\newtheorem{lemma}[theorem]{Lemma}}{}
\makeatother
\title[On lower bounds for hypergeometric tails]{On lower bounds for hypergeometric tails}
\author{Jianhang Ai, Christos Pelekis}
\date{Source: arXiv:2601.09485v2 (CC BY 4.0)}
\begin{document}
\maketitle

\noindent\emph{Source and licence.} On lower bounds for hypergeometric tails. Version arXiv:2601.09485v2 (CC BY 4.0), distributed under the Creative Commons Attribution 4.0 International licence, as recorded in the arXiv licence metadata for this exact version and shown on the abstract page https://arxiv.org/abs/2601.09485v2. Full rights notice: licenses/arxiv-2601.09485-CC-BY-4.0.txt.\par\smallskip

\noindent\textit{Editorial scope.} Counted unit: the Lemma whose source label is \texttt{mms1} in the article (source0.tex lines 672--678, source label \texttt{mms1}), delivered together with the article's own complete original proof of it (source0.tex lines 679--735, one \texttt{proof} environment of 57 lines). No mathematics is written, repaired or re-derived: statement and proof are the article's own text line for line. Required context retained uncounted: clean\_form (lines 278--286); main\_assumption1 (lines 665--667). Every definition, display and label of those lines is retained unchanged. Omitted from this package and not redistributed: 1--277, 287--664, 668--671, 736--998. Nothing inside a retained excerpt is omitted or altered: every retained body is delivered exactly as the article's lines are written, so no omission receipt of any kind accompanies this package. Citations: the retained excerpts carry no citation command at all, so no bibliography entry of the article is transcribed and no reference is redistributed. Reference adaptation: none was needed and none was made; every reference inside the delivered unit resolves inside it, so no unresolved reference or citation is produced. Printed slips kept literal: no printed word, symbol, hypothesis or number of the counted statement or of its proof was altered. Layout and class: the article's own class is \texttt{article} with packages including geometry, amssymb, amsmath, palatino, graphicx, babel, amscd, amsthm, color, enumerate, hyperref, comment; the wrapper is the lane's pinned \texttt{amsart} editorial wrapper at 10pt on A4, which restates the theorem-like environments the retained text needs and adds the article's own macro definitions that the retained text needs (none), with extra wrapper packages (bm, bbm, dsfont, xspace, cancel, mathtools). The delivered numbering is the unit's own. Assets: no retained span contains a figure environment, a graphics-inclusion command, a PDF-page inclusion, a table or a TikZ picture; no image file is copied into this package, the package declares no asset dependency and no image file is redistributed with it. Licence: version arXiv:2601.09485v2 of this article is distributed under the Creative Commons Attribution 4.0 International licence, as recorded in the arXiv licence metadata for this exact version and shown on the abstract page https://arxiv.org/abs/2601.09485v2. A full rights notice is delivered at licenses/arxiv-2601.09485-CC-BY-4.0.txt. Characters: every retained excerpt is pure ASCII, so the delivered engine, XeLaTeX with its default Latin Modern text font, reports no missing character.
\begin{lemma}\label{clean_form}
Let $H\sim\text{Hyp}(n,i,k)$, 
and fix an integer $m\in \{0,1,\ldots,\min\{i,k\}\}$. 
Then 
\begin{equation}\label{clean_form1}
\mathbb{P}(H \ge m) = \frac{\binom{k}{m}\cdot \binom{i}{m}}{\binom{n}{m}} \cdot \mathbb{E}\left(\frac{1}{\binom{Z_m+m}{m}} \right) \, ,
\end{equation}
where $Z_m\sim\text{Hyp}(n-m,i-m,k-m)$. 
\end{lemma}

\begin{equation}\label{main_assumption1}
k\cdot\frac{i}{n}\cdot \frac{n-i}{n}< 1 \quad \text{ and } \quad n\ge 8k \, .
\end{equation}

\begin{lemma}\label{mms1}
Suppose that~\eqref{main_assumption1} holds true, and that  $i<n/2$.  
Then  
\[
\mathbb{P}(H \ge ik/n) \, \ge \, \frac{k}{n} \, .
\]
\end{lemma}

\begin{proof}
Since $i<n/2$, it follows from~\eqref{main_assumption1}  that $\frac{ik}{n} < \frac{n}{n-i} < 2$, and thus $\frac{ik}{n}\in (0,2)$. 
We consider two cases. 

Suppose first that $\frac{ik}{n} \in (0,1]$.  
We may assume that $i\ge 2$; otherwise the result is clearly correct. 
Since $\frac{ik}{n} \le 1$ it holds $\mathbb{P}(H \ge ik/n) =\mathbb{P}(H \ge 1)$ 
and therefore it is enough to show that
\[
\mathbb{P}(H\ge 1)\ge \frac{k}{n} \, .
\] 
Now, Lemma~\ref{clean_form} implies that 
\[
\mathbb{P}(H\ge 1)=\frac{k}{n} \cdot \mathbb{E}\left(\frac{i}{Z_1 + 1}\right)   \, , 
\]
where $Z_1\sim\text{Hyp}(n-1,i-1,k-1)$.   The result follows upon observing that $i\ge Z_1+1$. 

Suppose now that $\frac{ik}{n} \in (1,2)$ and $i<\frac{n}{2}$. 
Observe that the assumption $n\ge 8k$ implies that $i> \frac{n}{k} \ge 8$, and so we may assume that $i\ge 9$. 
We may also assume that $k\ge 2$; otherwise the result is clearly true. 
Since $1<\frac{ik}{n} < 2$ it follows that 
$\mathbb{P}(H \ge \frac{ik}{n})=\mathbb{P}(H\ge 2)$. 
Moreover, it follows from Lemma~\ref{clean_form} that 
\[
\mathbb{P}(H\ge 2) = \frac{k(k-1)}{n(n-1)} \cdot i(i-1)\cdot \mathbb{E}\left(\frac{1}{(Z_2+2)(Z_2+1)}\right) \, ,
\]
where $Z_2\sim\text{Hyp}(n-2,i-2,k-2)$. 
Jensen's inequality now implies that 
\[
\mathbb{E}\left(\frac{1}{(Z_2+2)(Z_2+1)}\right) \ge 
\frac{1}{\mathbb{E}((Z_2+2)(Z_2+1))} = \frac{1}{\mathbb{E}(Z_2^2+3Z_2 +2)} \, .
\]
Since $Z_2\sim\text{Hyp}(n-2,i-2,k-2)$, we compute 
\[
\mathbb{E}(Z_2^2+3Z_2 +2) = A \cdot 
\frac{(n-i)(n-k)}{(n-2)(n-3)} + A^2 + 3A + 2 \, , \, \text{ where } \, A := \frac{(i-2)(k-2)}{n-2} \, .
\]
Hence, in order to prove the desired result, it is enough to prove that 
\[
\frac{k(k-1)}{n(n-1)} \cdot i(i-1)\cdot \frac{1}{\mathbb{E}(Z_2^2+3Z_2 +2)} \ge \frac{k}{n} \, ,
\]
which is equivalent to
\begin{equation}\label{ineq_mm1}
A \cdot 
\frac{(n-i)(n-k)}{(n-2)(n-3)} + A^2 + 3A + 2 \le i\cdot \frac{(i-1)(k-1)}{n-1} \, .
\end{equation}
Since $\frac{(n-i)(n-k)}{(n-2)(n-3)}<1$, \, $A\le B:=\frac{(i-1)(k-1)}{n-1}$ and $i\ge 9$, we conclude that~\eqref{ineq_mm1} will follow upon showing that 
\begin{equation}\label{ineq_mm2}
B^2 +4B + 2 \le 9B \quad \Longleftrightarrow\quad B^2-5B+2 \le 0 \, .
\end{equation}
Solving for $B$, we find that~\eqref{ineq_mm2} holds true for $B\in \left[\frac{1}{2}(5-\sqrt{17})\, ,\, \frac{1}{2}(5+\sqrt{17})\right]$; hence, since $B\le \frac{ik}{n}\le 2$, the proof will be complete after showing that $B > \frac{1}{2}(5-\sqrt{17})\approx 0.43844$. 
To this end, observe that the assumptions $i\ge 9$ and $i<\frac{n}{2}$ imply that 
\[
\frac{B}{\frac{ik}{n}}=\frac{\frac{(i-1)(k-1)}{n-1}}{\frac{ik}{n}} = \frac{n}{n-1} \cdot \left(1-\frac{1}{i}\right)\cdot \left(1-\frac{1}{k}\right) > \frac{8}{9}\cdot \left(1-\frac{i}{n}\right) >\frac{8}{9}\cdot\frac{1}{2} > 0.44 \, .
\]
Since  $\frac{ik}{n}>1$, it follows that $B > 0.44 \cdot \frac{ik}{n} > 0.44$,  and the result follows. 
\end{proof}

\end{document}
